Large-scale Internet outage detection combining BGP routing data, active
probing, and network-telescope (Internet background radiation) traffic is
the general measurement approach this paper builds on, developed and
operated as IODA~\citep{alcock2024ioda}. This paper adopts IODA's own
BGP-visibility convention (a route is counted as visible if seen by at
least half of the route collectors monitoring it) directly, rather than
re-deriving an independent threshold, for comparability
(Section~\ref{sec:methods}).

Iran's November 2019 nationwide blackout -- a short reference event that
falls inside this study's pre-shutdown baseline period (Section~\ref{sec:methods})
-- was previously characterized through network-telescope and
compromised-host measurement combined with
BGP reachability tracking,
finding cellular operators disconnected first and the bulk of the country
following over the next several
hours~\citep{mangino2021multidimensional}. This paper's H4 comparison
(Section~\ref{sec:h4}) revisits that same event at BGP-update-stream
resolution alongside three later windows, using a straggler-robust
interpercentile duration metric rather than full withdrawal range.

Two contemporaneous, independently authored studies cover parts of the
same 2025--26 shutdown period this paper measures, each with a distinct
methodology. \citet{aceto2026iran} cover the January 2026 event identified
in Section~\ref{sec:h1} using a combination of public BGP/RIS/RouteViews
data, IODA active probing, and network-telescope data -- the same core
measurement trio this paper uses, plus CDN/passive-flow data and OONI
censorship-measurement data this paper does not have. This paper
frames its own treatment of the January event as an independent
reproduction via measurement, adding finer BGP-level onset/offset timing to
prior public coverage instead of presenting the event as new
(Section~\ref{sec:h1}).
\citet{jahromi2026multiperspective} cover the 2026 blackout using passive
Censys scans and active TCP probing from external vantage points combined
with BGP analysis -- a measurement approach with no overlap in data
sources with this paper's own (which relies on BGP routing data and
IODA's active-probing and network-telescope signals) -- and report the shutdown enforced
via forwarding-plane null-routing with BGP announcements remaining stable,
plus structural exemptions for academic networks and a major CDN. Where
this paper's own findings intersect with theirs (Sections~\ref{sec:h1},
\ref{sec:h3}), the two independently derived measurements are pointed out
as corroborating rather than merged, since the underlying vantage points
and methods differ.

Beyond these Iran-specific studies, this paper's approach sits inside a
broader body of censorship-measurement work. Combining BGP control-plane
data with network-telescope (darknet) traffic to detect and characterize
country-wide outages -- the core combination this paper also uses --
originates with \citet{dainotti2013outages}'s analysis of the 2011 Egypt
and Libya shutdowns. \citet{bischof2023destination} extend that approach
to the first systematic, cross-country characterization of
government-ordered shutdowns versus spontaneous outages (155 countries
over four years), built on the same IODA dataset this paper draws
on~\citep{alcock2024ioda}; this paper's single-country, two-series
measurement design is a deep case study within the space that survey
characterizes at scale. \citet{wendzel2025survey} situate both within the
wider internet-censorship-measurement field more broadly, covering
measurement methodology and available datasets beyond outages and
shutdowns specifically.
